\begin{solapartat}
\begin{minipage}[t]{0.6\linewidth}
\vspace{0pt}
\punts{0,2} $F_m = ma_c$

\smallskip
\punts{0,2} $qvB = m\dfrac{v^2}{r} \Rightarrow v = \dfrac{qBr}{m}$

\smallskip
\punts{0,2} $v = \dfrac{1,60\times 10^{-19}\times 0,50\times 0,30}{1,70\times 10^{-27}} = 1,41\times 10^{7}\ \mathrm{m\,s^{-1}}$
\end{minipage}\hfill
\begin{minipage}[t]{0.07\linewidth}
\vspace{0pt}
\punts{0,4}
\end{minipage}
\begin{minipage}[t]{0.3\linewidth}
\vspace{0pt}
\figura[1]{sol_fig1.png}
\end{minipage}
\end{solapartat}

\begin{solapartat}
\punts{0,2} $T = \dfrac{2\pi r}{v}$

\smallskip
\punts{0,2} $T = \dfrac{2\pi\cdot 0,3}{1,41\times 10^{7}} = 1,33\times 10^{-7}\ \mathrm{s}$

\smallskip
\punts{0,2} $\omega = \dfrac{2\pi}{T} = 4,72\times 10^{7}\ \mathrm{rad\,s^{-1}}$

\smallskip
\punts{0,2} $E_c = \dfrac{1}{2}mv^2 = \dfrac{1}{2}1,70\times 10^{-27}\left(1,41\times 10^{7}\right)^2 = 1,69\times 10^{-13}\ \mathrm{J}$

\smallskip
\punts{0,2} Com que $F$ és perpendicular a la velocitat durant tot el moviment, no hi ha acceleració tangencial. El mòdul de la velocitat és constant (MCU). L'energia cinètica després d'una volta és la mateixa.
\end{solapartat}
